\documentclass[a4paper,11pt]{article}

% ----- Packages -----
\usepackage[margin=1.6cm]{geometry}
\usepackage{titlesec}
\usepackage{enumitem}
\usepackage{xcolor}
\definecolor{cvlink}{HTML}{174D70}
\usepackage[colorlinks=true,allcolors=cvlink]{hyperref}
\usepackage{setspace}
\usepackage{fontawesome5}
\usepackage{ifthen}

\setlength{\parindent}{0pt}

% ----- Section formatting -----
\titleformat{\section}{\Large\bfseries}{}{0em}{}[\titlerule]

% ----- Compact spacing -----
\setlist[itemize]{leftmargin=1.4em, itemsep=0.2em}

% =====================================================
% Resume Commands
% =====================================================

\newcommand{\cvLink}[2]{\href{#1}{\underline{#2}}}
\newcommand{\cvResource}[2]{[\cvLink{#1}{#2}]}

% #1 Title
% #2 Subtitle (institution/status/link text)
% #3 Date/location
\newcommand{\resumeEntry}[3]{
    \textbf{#1}\\
    \makebox[\textwidth]{
        #2
        \hfill
        #3
    }\\[-1.2em]
}

\newcommand{\simpleEntry}[2]{
    \textbf{#1}
    \hfill
    #2\\[-1.2em]
}

\newcommand{\resumeBullets}[1]{
    \begin{itemize}[topsep=0.2em, itemsep=0.1em, parsep=0em]
        #1
    \end{itemize}
    \vspace{0.3em}
}

\newcommand{\resumeBullet}[1]{
    \item #1
}


\begin{document}

% =====================================================
% Header
% =====================================================

\begin{center}
    {\huge \textbf{Tairan Wang}} \\[14pt]
    \normalsize
    \cvLink{mailto:tairan.wang.22@ucl.ac.uk}{\faEnvelope\ \ tairan.wang.22@ucl.ac.uk}
    \quad
    \cvLink{https://github.com/WangTairan}{\faGithub\ \ Github}
    \quad
    \cvLink{https://wangtairan.github.io}{\faGlobe\ \ Homepage}
\end{center}


% =====================================================
\section*{Education}

\textbf{University College London} \hfill London, UK \\
MEng Mathematical Computation, \textbf{First-Class Honours (78\%)}
\hfill Sep 2022 -- Jun 2026



% =====================================================
\section*{Research Experience}


\resumeEntry
{Back to the Future: Regressing Readability Features from LLMs}
{\cvResource{https://wangtairan.github.io/resources/publications/back-to-the-future-regressing-readability-features-from-llms.pdf}{PDF}\enspace
\cvResource{https://arxiv.org/abs/2610.04641}{arXiv}\enspace
(First Author)}
{UCL \quad Oct 2026}

\resumeBullets{
    \resumeBullet{Developed BTTF, an interpretable code readability model
    combining linear regression with traditional code metrics,
    embedding-based semantic organisation, and language-model predictability.}

    \resumeBullet{Evaluated on 1,100 samples covering all six human-rated
    code-readability datasets in our study, achieving \textbf{state-of-the-art
    (SOTA)} overall performance; further validated robustness through
    controlled code transformations.}
}


\resumeEntry
{When Specifications Conflict: A Symmetry-Based Framework for Measuring LLM Preferences}
{\cvResource{https://wangtairan.github.io/resources/publications/when-specifications-conflict.pdf}{PDF}\enspace
\cvResource{https://arxiv.org/abs/2607.28384}{arXiv}\enspace
(First Author)}
{UCL \quad Jul 2026}

\resumeBullets{
    \resumeBullet{Proposed a symmetry-based framework for measuring LLM
    preferences under conflicting specifications, evaluated on 550 executable
    conflicts across 11 mathematical function families, with extensions to
    Boolean algebra, code generation, and clinical rules.}
}



\resumeEntry
{Cognitive Modeling of Code Readability: Working Memory and Semantic Recoverability}
{\cvResource{https://wangtairan.github.io/resources/publications/cognitive-modeling-of-code-readability.pdf}{PDF}\enspace
MEng Dissertation}
{UCL \quad Jun 2026}

\resumeBullets{
    \resumeBullet{Developed CognaScore and PreFixScore to model working-memory
    load and semantic recoverability in code comprehension; evaluated both on
    function- and file-level human-rated readability datasets.}
}



\resumeEntry
{Meow-Omni 1: A Multimodal Large Language Model for Feline Ethology}
{\cvResource{https://wangtairan.github.io/resources/publications/meow-omni-1.pdf}{PDF}\enspace
\cvResource{https://neurips.cc/virtual/2026/poster/151616}{NeurIPS 2026}}
{UCL \quad May 2026}

\resumeBullets{
    \resumeBullet{Developed an open-source quad-modal large language
    model integrating video, audio, physiological time-series, and textual
    reasoning for feline intent recognition.}
}



\resumeEntry
{HoarePrompt: Structural Reasoning About Program Correctness in Natural Language}
{\cvResource{https://wangtairan.github.io/resources/publications/hoareprompt.pdf}{PDF}\enspace
\cvResource{https://conf.researchr.org/details/icse-2026/icse-2026-research-track/164/HoarePrompt-Structural-Reasoning-About-Program-Correctness-in-Natural-Language}{ICSE 2026}}
{UCL \& PKU \quad May 2024--Jan 2025}

\resumeBullets{
    \resumeBullet{Proposed a framework for verifying program correctness from
    natural-language specifications using Hoare triples, strongest
    postcondition calculus, and \(k\)-induction for loops.}
}



% =====================================================
\section*{Professional Experience}


\resumeEntry
{CTO, Alephedge}
{\cvResource{https://alephedge.com}{Website}\enspace
\cvResource{https://alephedge.com/zh/products/ordinal/}{Ordinal}}
{Hong Kong \quad Sep 2026 -- Present}

\resumeBullets{
    \resumeBullet{Lead the design and development of Ordinal, an AI agent
    combining general-purpose capabilities with financial research expertise,
    integrating tools for data analysis, factor evaluation, and strategy
    backtesting.}
}


\simpleEntry
{Co-Founder \& Core Developer, Alephedge}
{Hong Kong \quad May--Sep 2026}

\resumeBullets{
    \resumeBullet{Built and maintained the company intranet and public website.}
}



\begin{samepage}
\simpleEntry
{Summer Research Internship, UCL}
{London \quad Jun--Sep 2025}

\resumeBullets{
    \resumeBullet{Worked on code readability model.}
}
\end{samepage}



% =====================================================
\section*{Technical Skills}

\resumeBullets{
    \resumeBullet{\textbf{Programming:} Python, Java, C.}

    \resumeBullet{\textbf{AI \& Machine Learning:} Machine Learning,
    Deep Learning, Large Language Models (LLMs), multimodal learning, and
    AI agent development.}

    \resumeBullet{\textbf{Software Engineering:} Full-stack development,
    AI-assisted development workflows, and
    maintainable software engineering practices.}

    \resumeBullet{\textbf{Quantitative Computing:} Factor research,
    portfolio construction, risk management.}
}



% =====================================================
\section*{Academic Achievements}

\resumeBullets{
    \resumeBullet{Ranked first in multiple UCL Computer Science modules,
    including COMP0010 (Software Engineering, \textbf{1/150}), COMP0169
    (Machine Learning for Visual Computing, \textbf{1/130}), and COMP0024
    (Artificial Intelligence and Neural Computing, \textbf{1/50}).}
}



\end{document}
